\chapter{Парламент}

\chapter{Органи судової влади}

Судова влада в Україні реалізується через систему судів загальної юрисдикції та
Конституційний Суд. Побудова КСЗІ для органів судової влади вимагає врахування
специфічної ієрархічної структури судів, органів суддівського врядування та
самоврядування, а також установ, що забезпечують діяльність судів. Єдиною
інформаційно-комунікаційною платформою судової гілки влади є \textbf{ЄСІКС}~---
Єдина Судова Інформаційно-Комунікаційна Система, яка складається з 10~розділів
(2~томи технічних вимог) і охоплює весь спектр судочинних, адміністративних та
аналітичних процесів.

\newpage
\subsection*{Структурні підрозділи}

\begin{enumerate}
    \item Верховний Суд
    \begin{itemize}
        \item Велика Палата Верховного Суду
        \item Касаційний адміністративний суд
        \item Касаційний господарський суд
        \item Касаційний кримінальний суд
        \item Касаційний цивільний суд
    \end{itemize}
    \item Вищі спеціалізовані суди
    \begin{itemize}
        \item Вищий суд з питань інтелектуальної власності
        \item Вищий антикорупційний суд
        \item Спеціалізований окружний адміністративний суд
        \item Спеціалізований апеляційний адміністративний суд
    \end{itemize}
    \item Апеляційні суди
    \item Місцеві суди (загальні, господарські, окружні адміністративні)
    \item Органи суддівського врядування
    \begin{itemize}
        \item Вища рада правосуддя (ВРП)
        \item Вища кваліфікаційна комісія суддів України (ВККСУ)
    \end{itemize}
    \item Органи суддівського самоврядування
    \begin{itemize}
        \item Рада суддів України
        \item Громадська рада доброчесності (ГРД)
        \item Громадська рада міжнародних експертів (ГРМЕ)
    \end{itemize}
    \item Установи забезпечення діяльності судів
    \begin{itemize}
        \item Державна судова адміністрація (ДСА)
        \item Територіальні управління ДСА, допоміжні установи
        \item Національна школа суддів України
        \item Служба судової охорони
    \end{itemize}
\end{enumerate}

\newpage
\section{Ієрархія OID профілів безпеки судової системи}

Кожен орган та рівень судової системи отримує унікальний ідентифікатор об'єкта
(OID) в гілці \texttt{1.2.804.3.1.2}, зареєстрованій у національному дереві OID
України (\texttt{iso.member-body.ua}). Цей OID вбудовується в X.509 сертифікат як
розширення Certificate~Policy (\texttt{id-ce-certificatePolicies}), що дозволяє
апаратно ідентифікувати рівень безпеки та інституційну приналежність кожного
суб'єкта системи. Імплементація виконана в модулі \texttt{CA.CP} бібліотеки
\texttt{synrc/ca}.

\subsection{Базовий та галузевий профілі}

\begin{tabular}{ll}
\textbf{OID}                 & \textbf{Профіль} \\
\hline
\texttt{1.2.804.3.1.2.1}     & Базовий профіль безпеки (Level~1) \\
\texttt{1.2.804.3.1.2.2}     & Галузевий профіль безпеки (Level~2) \\
\texttt{1.2.804.3.1.2.3}     & Цільовий профіль безпеки ЦОД (Level~3) \\
\texttt{1.2.804.3.1.2.4}     & Цільовий профіль безпеки Судів (Level~3) \\
\end{tabular}

\vspace{0.5em}
Базовий профіль визначає мінімально необхідний набір контролів для обробки відкритої
та конфіденційної інформації в судовій системі. Галузевий профіль розширює базовий
набір контролів відповідно до специфіки обробки чутливої інформації в державних
реєстрах та судах, враховуючи вимоги високої доступності й цілісності даних.

Мінімальна кількість цільових профілів безпеки Level~3, необхідних для функціонування
будь-якого суду,~--- \textbf{два}:

\begin{enumerate}
    \item \textbf{Цільовий профіль ЦОД}~--- профіль безпеки центру обробки даних,
          який базується на галузевому профілі Міністерства цифрової трансформації
          України (DIGITAL UKRAINIAN GOVERNMENT INDUSTRY PROFILE). Визначає контролі
          фізичної безпеки, мережевої інфраструктури, резервного копіювання та
          катастрофостійкості серверного середовища, на якому розгортається ЄСІКС.
    \item \textbf{Цільовий профіль COURT}~--- базовий цільовий профіль для всіх
          судів. Визначає контролі прикладного рівня: управління доступом (ABAC),
          криптографічний захист судових документів (X.509, CMS), аудит дій
          користувачів, цілісність реєстрів та захист персональних даних учасників
          судових процесів.
\end{enumerate}

Кожен конкретний суд або орган правосуддя може додатково мати власний цільовий
профіль Level~3 або Level~4, який успадковує і розширює профіль COURT відповідно
до специфіки своєї юрисдикції.

\newpage
\subsection{Цільовий профіль вищих судів}

\begin{tabular}{ll}
\textbf{OID}                 & \textbf{Суд} \\
\hline
\texttt{1.2.804.3.1.2.4.1}     & Цільовий профіль вищих судів \\
\texttt{1.2.804.3.1.2.4.1.1}   & Велика Палата Верховного Суду \\
\texttt{1.2.804.3.1.2.4.1.2}   & Касаційний адміністративний суд \\
\texttt{1.2.804.3.1.2.4.1.3}   & Касаційний господарський суд \\
\texttt{1.2.804.3.1.2.4.1.4}   & Касаційний кримінальний суд \\
\texttt{1.2.804.3.1.2.4.1.5}   & Касаційний цивільний суд \\
\end{tabular}

\subsection{Цільовий профіль вищих спеціалізованих судів}

\begin{tabular}{ll}
\textbf{OID}                 & \textbf{Суд} \\
\hline
\texttt{1.2.804.3.1.2.4.2}     & Вищі спеціалізовані суди \\
\texttt{1.2.804.3.1.2.4.2.1}   & Вищий суд з питань інтелектуальної власності \\
\texttt{1.2.804.3.1.2.4.2.2}   & Вищий антикорупційний суд \\
\texttt{1.2.804.3.1.2.4.2.3}   & Спеціалізований окружний адміністративний суд \\
\texttt{1.2.804.3.1.2.4.2.4}   & Спеціалізований апеляційний адміністративний суд \\
\end{tabular}

\subsection{Цільовий профіль місцевих та апеляційних судів}

\begin{tabular}{ll}
\textbf{OID}                 & \textbf{Профіль} \\
\hline
\texttt{1.2.804.3.1.2.4.3}     & Цільовий профіль судів (Level~3) \\
\texttt{1.2.804.3.1.2.4.3.1}   & Місцеві суди (Level~4) \\
\texttt{1.2.804.3.1.2.4.3.2}   & Апеляційні суди (Level~4) \\
\end{tabular}

\vspace{0.5em}
Місцеві суди охоплюють загальні, господарські, окружні та адміністративні суди
першої інстанції. Кожна категорія успадковує контролі Level~3 та доповнює їх
специфічними вимогами до обробки судової документації відповідної юрисдикції.

\subsection{Органи та установи системи правосуддя}

\begin{tabular}{ll}
\textbf{OID}                 & \textbf{Орган / Установа} \\
\hline
\texttt{1.2.804.3.1.2.5}     & Органи та установи правосуддя \\
\texttt{1.2.804.3.1.2.5.1}   & ДСА~--- Державна судова адміністрація \\
\texttt{1.2.804.3.1.2.5.2}   & ТУ ДСА, допоміжні установи \\
\texttt{1.2.804.3.1.2.5.3}   & Рада суддів України \\
\texttt{1.2.804.3.1.2.5.4}   & ВРП~--- Вища рада правосуддя \\
\texttt{1.2.804.3.1.2.5.5}   & ВККСУ \\
\texttt{1.2.804.3.1.2.5.6}   & Національна школа суддів України \\
\texttt{1.2.804.3.1.2.5.7}   & ГРД та ГРМЕ~--- громадські ради \\
\texttt{1.2.804.3.1.2.5.8}   & Служба судової охорони \\
\end{tabular}

\newpage
\section{Таксономія компонент ЄСІКС}

Єдина Судова Інформаційно-Комунікаційна Система (ЄСІКС) поділяється на 10~функціональних
розділів, згрупованих у два томи технічних вимог. Нижче наведено повну таксономію з
кодифікацією субпродуктів.

\subsection{I. Інфраструктурні сервіси}
\begin{enumerate}
    \item Інфраструктура безпеки (X.509, EUDI)
    \item Рольова модель і система доступів (ABAC, IAM)
    \item Оркестрація процесів (BPMN)
    \item Словникова підсистема (HL7)
    \item Мета сервіси схеми даних (SCHEMA)
    \item Конструктор форм і процесів (X-FORMS)
\end{enumerate}

\subsection{II. Господарська діяльність}
\begin{enumerate}
    \item Бюджетне планування і контроль видатків (BUD)
    \item Фінансовий і бухгалтерський облік (FIN)
    \item Кадрова система (ACC)
    \item Публічний документообіг і сервісне обслуговування (CRM)
\end{enumerate}

\subsection{III. Інфраструктура документообігу}
\begin{enumerate}
    \item Шаблони документів (DOCX, XLSX)
    \item Генерація документів (PDF)
    \item Генерація бар- та штрих-кодів (QR)
    \item Розпізнавання Tesseract (OCR)
    \item Текстуально-голосові трансформації (S2T, T2S)
    \item Редактор справ консиліуму (CRDT)
\end{enumerate}

\subsection{IV. Реєстри}
\begin{enumerate}
    \item Виконавчі документи (EXEC)
    \item Судові рішення (ORDER)
    \item Правові позиції (POSITION)
    \item Дайджести (DIGEST)
    \item Фізичні особи (CITIZEN)
    \item Адвокати (ESQ)
    \item Уповноважені (POA)
    \item Кадри судової влади (JUDGE)
    \item Суб'єкти господарювання (BIZ)
    \item Зовнішні реєстри (TREMBITA)
\end{enumerate}

\subsection{V. Судопис}
\begin{enumerate}
    \item Справи (DEEDS)
    \item Провадження (PROSEC)
    \item Судові засідання і їх рішення (COURT)
    \item Авторозподіл справ (AUTOCASE)
    \item Реєстрово-інтеграційні сервіси (CART)
\end{enumerate}

\subsection{VI. Суддівство}
\begin{enumerate}
    \item Облік суддівських кадрів (ACC)
    \item Відбір суддів (FILTER)
    \item Кваліфікаційне тестування (QUALIFY)
    \item Досьє на суддю (PROFILE)
    \item Верифікація (VERIFY)
\end{enumerate}

\subsection{VII. Інтелектуальні сервіси}
\begin{enumerate}
    \item Короткий зміст (EXCERPT)
    \item Рекомендації (RECOMMEND)
    \item Практика (PRACTICE)
    \item Адвокат (ATTORNEY)
    \item Сутності (ENTITIES)
\end{enumerate}

\subsection{VIII. Бізнес, аналіз і звітність}
\begin{enumerate}
    \item Календарі та завдання (CAL, TASK)
    \item Платежі (PAYMENT)
    \item Архів (ARCHIVE)
    \item Коментування (COMMENTS)
    \item Конференц-зв'язок (VKZ)
    \item Журнали (JOURNALS)
\end{enumerate}

\subsection{IX. Інформаційно-пошукова система (court.gov.ua)}
\begin{enumerate}
    \item Контент (CMS)
    \item Форми (X-FORMS)
    \item Інтеграція (TREMBITA)
    \item Звернення (CITIZENS)
\end{enumerate}

\subsection{X. Кабінет (cabinet.court.gov.ua)}
\begin{enumerate}
    \item Заяви (APP)
    \item Листи (LETTER)
    \item Звернення (REQ)
    \item Довіреності (POA)
    \item Виконавчі документи (EXEC)
    \item Провадження (PRO)
    \item Судові рішення (ADJ)
\end{enumerate}

\newpage
\section{Політики проєкту ЄСІКС}

Політики є первинним і незмінним документом будь-якого проєкту. Вони визначають
правила до того, як формулюються вимоги чи обираються технології. Основні
принципи, що лежать в основі ЄСІКС:

\begin{enumerate}
    \item Політики є первинними документами; всі рішення приймаються на їх основі.
    \item Технічні вимоги формуються як похідні від політик.
    \item Технічне завдання~--- єдиний документ, де дозволяється вказувати конкретні продукти.
    \item Архітектура залишається агностичною до моменту затвердження ТЗ.
    \item Технічні вимоги містять перелік сутностей, процесів, прав і привілеїв.
    \item Технічне завдання містить BPMN діаграми, ABAC правила, словники, протоколи.
    \item Технічна документація визначається стандартами ISO і ДСТУ.
\end{enumerate}

\section{Вимоги до безпеки ЄСІКС}

Вимоги до інформаційної безпеки ЄСІКС повинні відповідати:

\begin{itemize}
    \item Закону України <<Про захист інформації в інформаційно-телекомунікаційних системах>>;
    \item Закону України <<Про захист персональних даних>>;
    \item Закону України <<Про основні засади забезпечення кібербезпеки України>>;
    \item Правилам забезпечення захисту інформації (Постанова КМУ від 29.03.06 №373);
    \item НД ТЗІ 3.6-006-24 (гармонізовано з NIST SP~800-53);
    \item Стандарту OWASP TOP~10 для безпеки веб-застосунків.
\end{itemize}

Додатково: засвідчення даних здійснюється виключно КЕП у форматі CAdES-X~Long;
використовуються засоби КЗІ з чинним позитивним експертним висновком Держспецзв'язку;
TLS версії не нижче 1.3; захист від SQL-ін'єкцій, XSS, Command Execution.

\section{Таксономія системних компонент}

Системні компоненти ЄСІКС класифікуються за функціональними категоріями:

\begin{itemize}
    \item \textbf{Data Storages:} MongoDB, SQL, S3, Cassandra, Riak, TiKV.
    \item \textbf{PubSub Buses:} Kafka, MQTT (ISO), AMQP (ISO), XMPP (ISO).
    \item \textbf{Application Services:} Anti Fraud, OTP, Auth, Integration Layer,
          Digital Signature, MCI, ETL.
    \item \textbf{Facade Management:} Traefik / KONG (OSI Layer 4--7).
    \item \textbf{Virtual Machines:} JVM, CLR, BEAM.
    \item \textbf{External Systems:} Трембіта (ДССУ, Дія, ПФУ, Мінюст).
\end{itemize}

\addtocontents{toc}{\protect\newpage}
\newpage
\chapter{Органи виконавчої влади}

Цей розділ описує структуру органів виконавчої влади (ОВВ), які
виступають ключовими об'єктами інформатизації.

\section{Міністерство освіти і науки України}

Цей підрозділ є статтею-дослідженням таксономії структури Міністерства освіти і науки як з погляду інформаційної автоматизованої системи, так і в розрізі соціальної структури державного управління. Як приклад, тут наведено цільову предметну модель, що є оптимізованою модифікацією чинної ієрархічної системи Міністерства освіти і науки України.

\subsection*{Структурні підрозділи}

\begin{enumerate}
    \item Патронатна служба
    \item Управління початкової школи
    \item Управління середньої школи
    \item Управління вищої школи
    \item Управління юстиції
    \begin{itemize}
        \item Департамент експертизи і сертифікації
        \item Інститут інтелектуальної власності
        \item Департамент кадрового забезпечення
        \item Департамент аудиту та внутрішніх розслідувань
        \item Департамент архівної справи (SCAN)
        \item Технічний департамент (CA)
        \item Департамент соціального і гуманітарного забезпечення
        \item Відділ кадрів (ACC)
        \item Юридичний департамент
        \item Департамент міжнародного співробітництва
    \end{itemize}
    \item Управління науково-дослідними інститутами (агенція)
    \item Національна академія наук (агенція)
    \begin{itemize}
        \item Інститут формальної математики
        \begin{itemize}
            \item Ректорат формальної філософії
            \item Ректорат чистої математики
            \item Ректорат прикладної математики
            \item Ректорат мовного забезпечення
            \item Ректорат теоретичної інформатики
        \end{itemize}
        \item Інститут формальної літератури
        \item Інститут музики, кіно й образотворчого мистецтва
        \begin{itemize}
            \item Національна консерваторія
            \item Національна академія мистецтв
            \item Національна кінематика
        \end{itemize}
        \item Інститут фізики і матеріалознавства
        \item Інститут геології і геохімії
        \item Інститут хімії і біології
        \item Інститут соціальних і гуманітарних наук
        \begin{itemize}
            \item Ректорат філософії
            \item Ректорат археології
            \item Ректорат національної історії
            \item Ректорат права
            \item Національна бібліотека
        \end{itemize}
    \end{itemize}
    \item Управління політиками і структурними підрозділами (агенція)
    \begin{itemize}
        \item Департамент комунікації (відділ кадрів)
        \item Департамент контролю виконання показників (CRM)
        \item Департамент планування переходу (аналіз процесів BPMN)
        \item Департамент трансформації (широкий спектр спеціалізацій)
    \end{itemize}
\end{enumerate}

\newpage
\section{Міністерство у справах ветеранів України}
\section{Міністерство охорони здоров'я України}

Таксономія Міністерства охорони здоров'я базується на послідовній інтеграції існуючих клінічних баз, системи громадського здоров'я та профільних дослідницьких інституцій у єдину організаційно-інформаційну мережу. Ця модель зорієнтована на оперативне забезпечення сталого розвитку галузі, медичної освіти й науки.

\subsection*{Структурні підрозділи}

\begin{enumerate}
    \item Патронатна служба
    \item Управління медичної освіти та науки
    \begin{itemize}
        \item Медичні університети й академії
        \item Науково-дослідні інститути (НДІ) системи МОЗ
    \end{itemize}
    \item Управління державного санітарного контролю та громадського здоров'я
    \begin{itemize}
        \item Центр громадського здоров'я (ЦГЗ)
        \item Департамент епідеміологічного нагляду
    \end{itemize}
    \item Департамент фармаконагляду і медичного ліцензування
    \item Управління юстиції
    \begin{itemize}
        \item Департамент медичної експертизи
        \item Департамент внутрішнього аудиту
        \item Департамент кадрового забезпечення та юридичної підтримки
    \end{itemize}
    \item Управління політиками та структурними підрозділами (агенція)
    \begin{itemize}
        \item Департамент цифрової трансформації (взаємодія з eHealth та ЕСОЗ)
        \item Департамент стратегічного планування
    \end{itemize}
\end{enumerate}

\newpage
\section{Міністерство внутрішніх справ України}

Архітектура МВС чітко розподіляє управлінські повноваження між своїм центральним апаратом та підзвітними йому центральними органами виконавчої влади (ЦОВВ). Це забезпечує ефективне реагування, незалежний внутрішній контроль і надання сервісних функцій населенню без дублювання повноважень.

\subsection*{Структурні підрозділи}

\begin{enumerate}
    \item Патронатна служба
    \item Головний департамент формування політик безпеки
    \item Управління координації ЦОВВ
    \begin{itemize}
        \item Національна поліція України (НПУ)
        \item Державна служба України з надзвичайних ситуацій (ДСНС)
        \item Державна прикордонна служба України (ДПСУ)
        \item Національна гвардія України (НГУ)
        \item Державна міграційна служба України (ДМС)
    \end{itemize}
    \item Головний сервісний центр (ГСЦ)
    \item Управління юстиції
    \begin{itemize}
        \item Департамент внутрішньої безпеки
        \item Департамент аудиту та службових розслідувань
    \end{itemize}
    \item Управління політиками та структурними підрозділами (агенція трансформації)
\end{enumerate}

\newpage
\section{Міністерство закордонних справ України}

Структурна модель МЗС відображає вимоги сучасного протоколу та міжнародного права. Вона оптимізована для швидкого обміну захищеною аналітичною інформацією з дипломатичними представництвами та для підтримки зовнішньополітичного консолідованого курсу.

\subsection*{Структурні підрозділи}

\begin{enumerate}
    \item Патронатна служба
    \item Управління дипломатичної освіти та аналітики
    \begin{itemize}
        \item Дипломатична академія
        \item Інститут актуальних міжнародних відносин
    \end{itemize}
    \item Територіальні управління
    \begin{itemize}
        \item Департаменти європейської, євроатлантичної, азійської та американської політики
    \end{itemize}
    \item Департамент міжнародних організацій (ООН, НАТО тощо)
    \item Управління юстиції
    \begin{itemize}
        \item Департамент міжнародного права
        \item Консульська служба
        \item Департамент кадрового аудиту та безпеки
    \end{itemize}
    \item Управління політиками та інституційного розвитку
\end{enumerate}

\newpage
\section{Міністерство оборони України}

Цей підрозділ є дослідженням таксономії структури Міністерства оборони України з погляду як інформаційної автоматизованої системи, так і соціальної структури в розрізі державного управління. Як приклад у статті розглядається конкретна структурна модель, що є адаптованою модифікацією чинної ієрархічної системи відомства. Ця структура підсилена повним спектром інституцій, необхідних для організації безперервного науково-освітнього та технологічно-виробничого циклу функціонування оборонного державного органу виконавчої влади.

Як модель гранулярності використано типову українську державну класифікацію (міністерство, управління, департамент, відділ, сектор). Оскільки ця робота зосереджена насамперед на логіці архітектури процесів, тут значною мірою надається перевага формальним моделям, які потребують мінімальних зусиль для своєї верифікації, моделювання та прогнозування. З огляду на те, що формалізація процесів безпосередньо стосується програмного забезпечення, у пригоді стають міжнародні стандарти телекомунікаційних протоколів. Їх сертифікація своєю чергою залучає університети, науково-дослідні та науково-виробничі інститути. Науково-виробничі інститути тут розглядаються у широкому сенсі --- як такі, що можуть функціонувати в симбіозі з комерційними об'єднаннями, міжнародними фундаціями тощо. Для адекватної підтримки автономної діяльності всі ці інституції повинні входити до єдиного арсеналу функціональних можливостей міністерства. Насамперед це стосується повноцінного університету четвертого рівня акредитації, навчальні програми якого мають перегукуватися (наприклад) з пропозиціями ключових кафедр Інституту математики НАНУ (див. Додаток 1). 

Фізичне розміщення базової інформаційної інфраструктури розгалуженої національної мережі міністерства та його підрозділів також передбачає автономне забезпечення (модель класу EDGE-офіс) із власними ресурсами охолодження, водо-електро-постачання та засобами посиленої охорони. Інформаційна політика формального моделювання допускає використання сучасних інструментальних мов розвитку, здатних підтримувати сувору машинну верифікацію на рівні комплексів ТЗІ Г7 (з можливостями повної математичної доводжуваності), покладаючись здебільшого на надійні телекомунікаційні протоколи та міжнародні стандарти серії ISO (див. Додаток 5).

\subsection*{Мотивація}

Основна мотивація цієї роботи полягає у висвітленні таксономії Міністерства оборони України через призму структурної оптимізації, максимальної інституційної автономності та життєстійкості (sustainability), а також здатності до самовідтворення й ефективної підтримки операційного життєвого циклу міністерства в умовах динамічних викликів.

\subsection*{Структурні підрозділи}

\begin{enumerate}
    \item Патронатна служба
    \item Управління освіти і науки (агенція)
    \begin{itemize}
        \item Університет четвертого рівня акредитації (див. Додаток 1)
        \item Науково-дослідні інститути (НДІ)
        \item Науково-виробничі інститути
    \end{itemize}
    \item Управління медицини (агенція)
    \begin{itemize}
        \item Клінічні наукові дослідження та лабораторії, НДІ профілактичної та військової медицини (MED)
        \item Заклади реабілітації («Пуща-Водиця», «Трускавецький», «Хмільник»)
        \item Клінічні та мобільні (4) лікарні, військові госпіталі (14)
        \item Медичні служби (5 родів військ), центри тактичної медицини, Командування Медичних сил
        \item Інститут медицини (на базі ЗДМУ, ХНМУ, ЛНМУ, ТНМУ)
    \end{itemize}
    \item Виробничо-промислове управління (агенція)
    \begin{itemize}
        \item Конструкторські бюро (КБ), концерн «Укроборонпром», ДАХК «Артем»
        \item Департамент економічного моделювання та планування
        \item Департамент ресурсного забезпечення (SCM, TMS, WMS)
        \item Департамент бюджетування і закупівель (FIN)
        \item Департамент будівництва, архітектури та масштабування ліній виробництва
        \item Телекомунікаційний департамент інформаційних систем
        \begin{itemize}
            \item Відділ систем урядування
            \item Відділ облікових систем
            \item Відділ телекомунікаційних систем
            \item Відділ безпекових протоколів інтернет-зв'язку
        \end{itemize}
        \item Департамент авіації, авіоніки, аеронавтики та безпілотних літальних апаратів (БПЛА) і їхніх систем
        \begin{itemize}
            \item Відділ авіації
            \item Відділ авіоніки
            \item Відділ аеронавтики
            \item Відділ безпілотних систем
        \end{itemize}
        \item Департамент машинобудування (тестування і проектування місцевості, SolidWorks)
        \item Департамент суднобудування та військово-морської техніки
    \end{itemize}
    \item Управління юстиції
    \begin{itemize}
        \item Відділ експертиз і сертифікації (ISO/IETF)
        \item Департамент архівної справи (SCAN)
        \item Департамент аудиту та внутрішніх розслідувань
        \item Технічний департамент (CA)
        \item Департамент соціального і гуманітарного забезпечення
        \item Відділ кадрів (ACC)
        \item Юридичний департамент
        \item Департамент міжнародного співробітництва
    \end{itemize}
    \item Управління розвідки
    \item Головне управління позиційною політикою та стратегією
    \begin{itemize}
        \item Мобілізаційний департамент
        \item Департамент військово-навчальних програм
        \item Департамент координації сил та засобів оборони (CRM, див. Додаток 2)
        \begin{itemize}
            \item Командування Сухопутних військ
            \item Командування Повітряних сил
            \item Командування Військово-Морських сил
            \item Командування спеціальних та допоміжних сил (ССО, Морська піхота, РЕБ, РХБЗ, ТРО)
            \item Командування кібербезпеки та ДШВ
        \end{itemize}
        \item Департамент контролю й оперативного управління DFR (див. Додаток 4)
        \item Економічний департамент
        \begin{itemize}
            \item Відділ матеріально-ресурсного забезпечення (WMS)
            \item Відділ бюджетування і закупівель
        \end{itemize}
    \end{itemize}
    \item Управління політиками та структурними підрозділами (агенція трансформації)
    \begin{itemize}
        \item Департамент комунікації та кадрової політики
        \item Департамент контролю виконання ключових показників (CRM)
        \item Департамент планування структурного переходу (аналіз процесів BPMN)
        \item Департамент цифрової та інституційної трансформації
    \end{itemize}
\end{enumerate}

\subsection*{Принципи}

Одним із головних принципів, закладених у фундамент Міністерства оборони (МО), є створення повноцінної екосистеми вищої профільної освіти, яка повинна функціонувати відповідно до найвищих міжнародних стандартів. Для задоволення фахових потреб та розвитку працівників і військовослужбовців на всіх рівнях цієї таксономії в основу організації роботи покладено безперервне функціонування відомчого університету четвертого рівня акредитації, чиї спеціальності мають повністю покривати наукові та технологічні запити самого міністерства.

Іншим універсальним і життєво необхідним принципом є принцип чіткого розподілу влади, з якого випливає наявність трьох ключових макрокорпусів МО.
\begin{enumerate}
    \item \textbf{Політичний корпус} (до якого належать головне управління позиційною політикою, управління політиками та переходом, а також патронатна служба). Він передбачає виокремлення спеціальної агенції трансформації, яка ініціює створення та виконує структурну апробацію інституцій міністерства. До цього корпусу також входить головне управління, що визначає стратегію застосування Збройних Сил України та відповідних засобів оборони.
    \item \textbf{Виконавчий корпус} (управління освіти і науки, медичне управління та надважливе виробничо-промислове управління). Цей корпус гарантує сталий розвиток високоінтелектуальних та ресурсоємних структурних підрозділів і виводить їх під автономний контроль профільних агенцій, що дозволяє їм гнучкіше взаємодіяти із зовнішніми технологічними й медичними екосистемами.
    \item \textbf{Судовий і правоохоронний корпус} (управління юстиції, галузевий трибунал). Це окремий, функціонально незалежний структурний блок, який реалізує процеси внутрішнього аудиту, комплаєнсу та службових розслідувань всередині соціоінформаційних систем оборонного відомства.
\end{enumerate}

\newpage
\subsection*{Самоорганізація}

Для систематизації процесу організаційного будівництва варто виділити ключові структурні етапи та засади:

\begin{itemize}
    \item Розподіл влади та мінімізація адміністративної таксономії;
    \item Нормалізація формальних процесів;
    \item Незалежний аудит міністерства та процес інституційної трансформації;
    \item Проєктування архітектури структурних підрозділів;
    \item Апробація результатів впровадження і циклічність процесу оновлення.
\end{itemize}

Розподіл влади традиційно залишається головним принципом ефективного державного управління. Контролюючі та ревізійні органи повинні спеціалізуватися винятково на перевірках і розслідуваннях, їхня структура має бути операційно виокремленою. Розвиток політик, регуляція бізнес-процесів управління та збір реквізитної інформації мають перебувати в компетенції політичного блоку, тоді як виконання завдань (наука, освіта, матеріальне виробництво та забезпечення) здійснюють суто виконавчі гілки відомства.

Вимоги до документування процесів і збереження даних повинні формуватися на найвищому технічному рівні. Це означає використання еквіваріантної семантики, спрощення архітектури процесів до нормальних форм, максимальну мінімізацію заплутаних горизонтальних зв'язків. Усі вони мають цілковито відповідати актуальним вимогам постанови КМУ №55. Водночас має бути організовано прозорий процес історіографії та цифрового аудиту діловодства й виробничих ланцюгів відповідно до суворого стандарту управління якістю ISO 19510 (що є важливим для сумісності з аудитом за стандартами НАТО). Запуск автоматизованого діловодства передбачається впроваджувати поступово і гранулярно: починаючи від ключових політико-формуючих центрів та їхніх найменш ресурсоємних проєктів, і далі розширюючи охоплення на всі інші департаменти та управління згідно з затвердженою магістральною стратегією.

\subsubsection{Патронатна служба}

Забезпечує сприяння реалізації стратегічних та політичних цілей Міністра оборони, безпосереднє персональне консультування, а також організаційне, інформаційно-публічне й експертно-аналітичне забезпечення його діяльності.

\subsubsection{Управління освіти й науки (агенція)}

Сучасний науковий потенціал, крім безпосередньо освітнього процесу, активно виділяє як науково-дослідний, так і науково-виробничий напрями. Сфери досліджень повністю диктуються актуальними потребами бойових або організаційних департаментів Міністерства оборони. Ці високотехнологічні компанії й лабораторії повинні знаходитися, як і підпорядкований університет, в єдиній орбіті стратегічного впливу МОУ.

\subsubsection{Виробничо-промислове управління (агенція)}

Концепція передбачає повне дзеркальне дублювання сфер безпосереднього виробництва засобів оборони відповідно до профілю керуючих департаментів, що формують політики ресурсного забезпечення. Оскільки військова промисловість консолідує базові ресурси оборонної індустрії, є методично доцільним тримати всю картину планування та виробництва під єдиною ієрархічною парасолькою міністерства. Наявні сьогодні виробничі хаби й об'єднання, такі як АТ «УОП» (колишній «Укроборонпром») та ДАХК «Артем», пропонується інтегрувати в систему як потужні зовнішні конструкторські бюро в єдиному переліку з профільними приватними комерційними підприємствами, з якими планується вибудовувати гнучку проєктну співпрацю.

\subsubsection{Управління юстиції}

Управління юстиції як незалежна гілка бере на себе функцію неформального обліку сил (розширений відділ кадрів та прав) і юридичного фіксування протокольних дій усередині системи. Далі вони систематично аналізуються антикорупційним відділом та департаментом внутрішніх розслідувань і аудиту. Саме тут зосереджено ключовий штат юридичної експертизи, що обслуговує всю масштабну ієрархію міністерства, а також координує правову взаємодію із міжнародними зовнішніми партнерами, такими як блок НАТО. Також під парасолькою цього управління доцільно тримати технічний департамент як автономний центр емісії криптографічних сертифікатів і ключів управління для всього штату працівників.

\subsubsection{Головне управління позиційною політикою та стратегією}

Це сучасне осмислення спрощених і водночас глобально стратегічних функцій Генерального штабу, який безпосередньо взаємодіє з наявними резервами Сил та засобів оборони. Загальна тактична і стратегічна консоль управління підпорядковує всю логіку збройної місії веденню магістральної позиційної політики: від ефективного розгортання сил та засобів до просунутого дата-орієнтованого прогнозування еволюції подій. Відповідно, правильна позиційна політика завжди спирається на максимально достовірний локальний і глобальний облік наявних ресурсів.

\subsubsection{Управління політиками та структурними підрозділами (агенція)}

Це формальний мозковий центр адміністративного врядування та інституційного переходу від застарілої структури до будь-якої наперед заданої оптимальної моделі (наприклад, тієї, що описана в цьому документі). Управління займається первинним аудитом чинних бізнес-процесів (їх юридичним і процедурним розтином) та проєктуванням їхньої поетапної трансформації в нові урядові й облікові цифрові системи. Свою діяльність структура тісно координує з телекомунікаційним департаментом, здійснюючи проєктний менеджмент: контролює процес впровадження, проводить апробацію систем і подальше їх калібрування.

Процес системної трансформації раціонально поділено на такі категорії (зворотно пропорційно до швидкоплинності впровадження):
\begin{enumerate}
    \item Оцифровування (діджиталізація) базових існуючих процесів без кардинальної модифікації їхньої логіки (наприклад, класичне діловодство).
    \item Створення архітектури нових функціональних процесів та надійного технологічного забезпечення для «вакантних» або новоутворених інноваційних екосистем, таких як виробничо-промислове і технологічне управління.
    \item Планування та стратегічний розвиток нових управлінь (це найдовший за часом процес, що охоплює розробку навчальних програм, побудову наукомістких продуктів та їх всебічну підтримку).
\end{enumerate}

Кожна фаза цього структурного метапроцесу вимагає побудови чіткого життєвого циклу проєкту, включно з продуктом, який лежить у його технологічній основі. До цього циклу обов'язково входить повний комплект супровідної документації: технічне завдання, ескізний проєкт, робочий технічний проєкт і детальна технологічна документація. Управління політиками природно розподіляє сфери компетенції з телекомунікаційним підрозділом і підтримує єдині стандарти документування процесів.

Ключові процеси й завдання управління переходом можна узагальнити так: 1) архівна систематизація процесного виробництва; 2) розробка гнучких протоколів запуску нових структурних підрозділів; 3) розробка сучасних процесів захищеного документообігу; 4) розробка процесів інтегрованого управління та субординаційної координації; 5) створення стандартів технологічних процесів на основі європейських стандартів сімейства ISO 9001.

\subsection*{Висновки та результати}

У статті представлена розширена і запропонована до впровадження таксономія оптимізованого і нормалізованого міністерства. Наша модель створена з урахуванням передового міжнародного та корпоративного досвіду організації структур, підкреслює чіткі принципи прозорого розподілу влади та організації як глобального (стратегічного), так і локального (операційного) контекстів для відомств. Цілі стають легко досяжними, адже організація функцій здійснюється в найпряміший та найефективніший спосіб, різко мінімізуючи складність та кількість горизонтальних протоколів взаємодії всередині самої системи.

У додатках до статті наведено таксономію навчальних та наукових програм для залученого університету, а також детальну ієрархію телекомунікаційного департаменту, серед основних функцій якого є розробка та практичне впровадження новітніх IТ-систем ЗСУ і МО. У процесі цього первинного концептуального моделювання ми спробували охопити всі ключові органи системи аж до найвищого рівня гранулярності — департаментів, відділів та секторів, вибудовуючи первинні індекси і їхні внутрішні протоколи взаємодії.

\subsection*{Бібліографія}

\begin{itemize}
    \item ДСТУ 2732-94 Діловодство й архівна справа. Терміни та визначення.
    \item Наказ МОУ від 26.05.2014 №333 «Інструкція з організації обліку особового складу».
    \item Наказ МВС від 21.11.2017 №608 «Порядок проведення службового розслідування».
    \item Постанова КМУ від 26.07.2018 №370 «Про затвердження Інструкції з діловодства».
    \item Наказ МОУ від 07.04.2017 №124 «Інструкція з діловодства в Міністерстві оборони України».
    \item Постанова КМУ від 07.10.2015 №393 «Положення про юридичну службу».
    \item Наказ МОУ від 29.11.2018 №604 «Інструкція з надання доповідей і донесень про події, кримінальні правопорушення, військові адміністративні правопорушення та адміністративні правопорушення, пов'язані з корупцією, порушення військової дисципліни...».
\end{itemize}

\newpage
\section{Міністерство юстиції України}

Мін'юст у цій архітектурі відіграє роль не лише галузевого регулятора правовідносин, але й центрального верифікаційного вузла для інфраструктурних процесів інших ОВВ. Міністерство виступає головним арбітром у реєстраційних діях та забезпечує дотримання законності на технічному та організаційному рівнях.

\subsection*{Структурні підрозділи}

\begin{enumerate}
    \item Патронатна служба
    \item Департамент державної реєстрації та нотаріату
    \item Департамент виконання кримінальних покарань
    \item Департамент судової роботи та експертного забезпечення
    \item Власний орган юстиції та внутрішнього контролю (Генеральна інспекція)
    \begin{itemize}
        \item Департамент внутрішніх розслідувань і комплаєнсу
        \item Департамент аудиту реєстрів
    \end{itemize}
    \item Управління цифровізації та структурних підрозділів
    \begin{itemize}
        \item Адміністрування державних реєстрів юстиції
        \item Забезпечення кіберстандарту (спільно з ЦОВВ)
    \end{itemize}
\end{enumerate}

\newpage
\section{Міністерство фінансів України}

Організаційна модель Міністерства фінансів відповідає за макрофінансову стабільність, податкову політику та бюджетування всіх державних ресурсів. Для ефективного управління вона суворо розмежовує блок аналітики, блок адміністрування надходжень та блок контролю.

\subsection*{Структурні підрозділи}

\begin{enumerate}
    \item Патронатна служба
    \item Департамент державного бюджету та фінансового планування
    \item Департамент податкової та митної політики
    \item Департамент управління державним боргом
    \item Управління юстиції
    \begin{itemize}
        \item Департамент фінансового аудиту
        \item Служба державного фінансового моніторингу
        \item Юридичний департамент
    \end{itemize}
    \item Управління політиками та структурними підрозділами (інтеграційне планування з Казначейством)
\end{enumerate}

\newpage
\section{Міністерство економіки України}

Архітектура цього міністерства сконцентрована на прогнозуванні макропоказників, регулюванні промислового та аграрного секторів, а також на питаннях захисту вільної конкуренції та залучення інвестицій.

\subsection*{Структурні підрозділи}

\begin{enumerate}
    \item Патронатна служба
    \item Департамент макроекономічного прогнозування та аналітики
    \item Департамент розвитку реального сектору економіки
    \item Департамент регуляції зовнішньоекономічної діяльності
    \item Управління юстиції
    \begin{itemize}
        \item Департамент антимонопольного контролю та нагляду
        \item Департамент внутрішнього аудиту
    \end{itemize}
    \item Управління політиками та структурними підрозділами
    \begin{itemize}
        \item Департамент методології закупівель (Prozorro)
    \end{itemize}
\end{enumerate}

\newpage
\section{Міністерство енергетики України}

В основу структурної моделі Міністерства енергетики покладено розподіл за ключовими сировинними та енергогенеруючими напрямами, з виділенням сильного блоку екологічної безпеки та стратегічної синхронізації з європейськими енергомережами.

\subsection*{Структурні підрозділи}

\begin{enumerate}
    \item Патронатна служба
    \item Департамент ядерної енергетики та атомно-промислового комплексу
    \item Департамент нафтогазового комплексу
    \item Департамент електроенергетичного комплексу та зеленої енергії
    \item Управління юстиції
    \begin{itemize}
        \item Інспекція з енергетичного нагляду
        \item Департамент внутрішнього аудиту та комплаєнсу
    \end{itemize}
    \item Управління політиками та структурними підрозділами (агенція переходу до ENTSO-E)
\end{enumerate}

\newpage
\section{Міністерство соціальної політики України}

Структура цього органу сфокусована на соціальному захисті, регулюванні ринку праці та пенсійному забезпеченні. Вона спирається на Єдину інформаційну систему соціальної сфери (ЄІССС) як інструмент діджиталізації соціальних послуг.

\subsection*{Структурні підрозділи}

\begin{enumerate}
    \item Патронатна служба
    \item Департамент пенсійного забезпечення та соціального страхування
    \item Департамент захисту прав дітей та піклування
    \item Департамент праці та зайнятості
    \item Управління юстиції (внутрішній контроль, координація інспекції праці)
    \item Управління політиками та структурними підрозділами (Трансформація ЄІССС)
\end{enumerate}

\newpage
\section{Міністерство регіональної політики України}

Таксономія спрямована на ефективне управління магістральними інфраструктурними проєктами, процесами післявоєнного відновлення та завершенням комплексної децентралізації.

\subsection*{Структурні підрозділи}

\begin{enumerate}
    \item Патронатна служба
    \item Департамент інфраструктури та житлово-комунального господарства
    \item Департамент регіонального розвитку та децентралізації
    \item Департамент містобудування, урбаністики та архітектури
    \item Управління юстиції
    \begin{itemize}
        \item Аудит будівельних і транспортних інспекцій
        \item Внутрішня безпека
    \end{itemize}
    \item Управління політиками та інституційного відновлення (ЄДЕССБ, логістичні системи)
\end{enumerate}

\newpage
\section{Міністерство цифрової трансформації України}

Мінцифра виступає головним ідеолог-архітектором всієї урядової метасистеми електронного врядування. Структура побудована як гнучка ІТ-агенція, поділена за продуктовими напрямами (додаток Дія, ШСД) та проєктними екосистемами (Дія.City).

\subsection*{Структурні підрозділи}

\begin{enumerate}
    \item Патронатна служба
    \item Департамент розвитку електронних послуг (екосистема «Дія»)
    \item Департамент розвитку ІТ-індустрії (Спеціальний режим «Дія.City»)
    \item Департамент розвитку цифрової інфраструктури (широкосмуговий доступ)
    \item Управління юстиції
    \begin{itemize}
        \item Департамент захисту персональних даних та криптографічного супроводу
        \item Департамент внутрішнього ІТ-аудиту
    \end{itemize}
    \item Департамент системної інституційної трансформації ОВВ (методична координація ІТ-директорів усіх міністерств)
\end{enumerate}

\newpage
\section{Міністерство захисту довкілля та природних ресурсів України}

Архітектура цього відомства забезпечує контроль за раціональним використанням природних ресурсів і впровадженням політик сталого розвитку згідно з екологічними директивами ЄС.

\subsection*{Структурні підрозділи}

\begin{enumerate}
    \item Патронатна служба
    \item Департамент управління відходами та екологічної безпеки
    \item Департамент лісового та водного господарства
    \item Департамент природно-заповідного фонду та надрокористування
    \item Управління юстиції
    \begin{itemize}
        \item Державна екологічна інспекція (наглядовий орган)
        \item Департамент внутрішнього аудиту
    \end{itemize}
    \item Управління політиками та структурними підрозділами (впровадження екологічних реєстрів)
\end{enumerate}
